\section{Preempção e Escalonamento}

Preempção é o mecanismo pelo qual o sistema operacional interrompe temporariamente a execução de uma aplicação — neste contexto, podendo ser uma \gls{vm} — para recuperar o
controle de um núcleo da CPU. Esse mecanismo é fundamental para a implementação do escalonamento, permitindo que múltiplas aplicações compartilhem um mesmo processador de forma eficiente.

Conforme discutido na seção de máquinas virtuais, cada \gls{vm} é representada internamente por uma \textit{Thread} no \gls{so} hospedeiro.
Dessa forma, o escalonamento de máquinas virtuais pode ser tratado como um caso particular do escalonamento de \textit{Threads}.

Nesse contexto, o processo de escalonamento do hipervisor pode ocorrer de duas formas: voluntário ou involuntário.

\subsection{Escalonamento Involuntário}

O escalonamento involuntário ocorre independentemente das ações da \textit{Thread} ou da \gls{vm} em execução. Nesse caso, o temporizador de hardware gera interrupções periódicas
que forçam a preempção da tarefa em execução, transferindo o controle para o kernel do hipervisor.
Ao receber essa interrupção, o escalonador é acionado e avalia qual \textit{Thread} deverá utilizar o núcleo no próximo intervalo de tempo. Caso seja necessário substituir a
\textit{Thread} atualmente em execução, o sistema realiza uma troca de contexto entre a \textit{Thread} interrompida e a nova \textit{Thread} selecionada para execução.

\subsection{Escalonamento Voluntário}

O escalonamento voluntário ocorre quando uma \textit{Thread} ou uma \gls{vm} passa a aguardar a disponibilidade de algum recurso, deixando voluntariamente o processador.
Durante esse período de inatividade, outras Threads podem ser executadas até que a entidade bloqueada esteja novamente apta a prosseguir. Esse tipo de escalonamento pode
ocorrer em diferentes partes do sistema, sendo um dos casos mais comuns aquele em que uma \gls{vm} necessita realizar uma operação de \gls{io} e precisa aguardar sua vez para acessar um
controlador de hardware.

Outra forma de escalonamento voluntário ocorre quando a \gls{vm} em execução entra em estado de idle. Nesse cenário, o hipervisor deve identificar essa condição e verificar a
existência de outras \glspl{vm} prontas para execução. Apenas na ausência de outras demandas, o sistema permite que a CPU entre efetivamente em estado de idle.

Entretanto, a detecção do momento em que uma \gls{vm} passa para o estado de idle não é trivial, e exige suporte específico da arquitetura para que essa transição seja
identificada.






